% !TITLE 如梦令·昨夜雨疏风骤
% !AUTHOR 李清照
% !DYNASTY 两宋
% !GENRE 词
% !ORDER 13
% !EDITION 《全宋词》，唐圭璋编，中华书局
% !EDITION_STATUS pending

\begin{poem}
昨夜雨疏风骤\textsuperscript{1}，浓睡不消残酒\textsuperscript{2}。
试问卷帘人\textsuperscript{3}，却道海棠依旧\textsuperscript{4}。
知否，知否？应是绿肥红瘦\textsuperscript{5}。
\end{poem}

\begin{annotations}
\notetext{1}{雨疏风骤：雨点稀疏但风势猛烈。疏，稀；骤，急。}
\notetext{2}{浓睡不消残酒：沉睡了一夜，醒来仍有酒意。浓睡，沉睡。}
\notetext{3}{卷帘人：正在卷起窗帘的侍女。一夜风雨后，词人最关心的是窗外花木是否安好。}
\notetext{4}{海棠依旧：海棠花（还和以前一样）。侍女随口一答，没有仔细看——她不懂词人的心事。}
\notetext{5}{绿肥红瘦：绿叶更加茂盛（肥），红花却凋零了（瘦）。“肥”和“瘦”本是形容人的体态，用在这里赋予了植物生命感。这是李清照最著名的创造。}
\end{annotations}

\begin{appreciation}
《如梦令》是李清照少女时代的作品，也是她留名词史的第一支曲子。全篇只有三十三个字，一问一答，像一出只演一瞬的独幕剧。三十三个字里藏着一个少女一整个早晨的心事。

“昨夜雨疏风骤”，雨点稀疏，风势却急。这样的雨打不湿衣裳，却能把满树的花在一夜之间摇落。“疏”与“骤”两个字并置，下笔便有分寸：她知道这雨的分量，也预感到花的命运。“浓睡不消残酒”，一夜沉睡，醒来酒意仍在。她为什么喝酒，词里没有说；这份惜花的敏感与伤春的怅惘大约是同一件事情的两端。

“试问卷帘人”，一个“试”字最见心思。想问，又怕听到不想听的答案，于是先在心里掂量几遍才敢开口。帘外的侍女随口一句“海棠依旧”，答得干脆，也答得敷衍，她多半没有细看，也读不懂主人为何这样在乎几枝花的开落。

于是词人急了：“知否，知否？应是绿肥红瘦。”叠问两遍，是纠正侍女，也是说服自己。一夜风雨过后，叶子喝饱了水，越发肥润；花经不住吹打，只剩下几瓣，越发清瘦。“肥”与“瘦”本是形容人的字眼，移来写花写叶，草木忽然有了肌肤与骨相。前人写花多半只写它开与谢的迟速浓淡，把花写成会胖会瘦的活物确是李清照独一份的本事。后人称她为“李三瘦”，其中一瘦便是这句“绿肥红瘦”；李清照另有一首《醉花阴》写“人比黄花瘦”，与此正是一脉相承。

这句也不是凭空而来。唐人韩偓《懒起》有“昨夜三更雨，临明一阵寒。海棠花在否？侧卧卷帘看”，意思相去不远，只是问得平直。李清照把前人一句家常的问话化成了“绿肥红瘦”四个字，一化便有了魂。

少年的愁，不过一场夜雨、几枝海棠。可就是这点愁，被她写得没有一滴眼泪，却比哭还要急。
\end{appreciation}
